\inicializarSeccion{Preguntas susceptibles de análisis inferencial}

Ya tenemos nuestra variable y nuestras preguntas definidas. Sin embargo, con estas preguntas realmente no podemos hacer análisis estadistico. Debido a que estas solo se limitan a describir un predicado; hay que reformularlas para que puedan ser analizadas con métodos estadísticos inferenciales. De esta forma, nosotros podemos realizar inferencias sobre la población y realizar pruebas de hipótesis, delimitaciones de confianza y otras técnicas estadísticas.

Para esto, se deben reformular las preguntas para que puedan ser analizadas con métodos estadísticos inferenciales:

\begin{itemize}
    \item ¿El nivel promedio de oxígeno disuelto en nuestra región cumple con el umbral mínimo óptimo establecido por las normativas ambientales (5 ml/L) necesario para mantener la salud del ecosistema?
    \item ¿De qué manera las variaciones e incrementos en la temperatura y la salinidad del agua están alterando significativamente los niveles de oxígeno disuelto (\texttt{mean\_O2ml\_L}), y cómo pone este déficit en riesgo el equilibrio de la biodiversidad marina?
    \item ¿Existe una relación significativa entre los niveles de oxígeno disuelto (\texttt{mean\_O2ml\_L}) y la biodiversidad marina en nuestra región?
\end{itemize}

Estas preguntas permiten realizar análisis estadísticos inferenciales para obtener conclusiones sobre la población y validar hipótesis relacionadas con la salud del ecosistema.